\section{其他并行策略简介}

\subsection{上下文并行 / Ring Attention}

\begin{figure}[H]
\centering
\includegraphics[width=0.95\textwidth]{slides-images/slide-050.jpg}
\caption{上下文并行（Ring Attention）示意\protect\footnotemark}
\end{figure}
\footnotetext{Slides 第51页。}

上下文并行（Context Parallelism）专门用于\textbf{超长序列}的注意力计算。核心思想：

\begin{itemize}
    \item 每个 GPU 负责不同位置的 Query
    \item Key 和 Value 以环形方式在 GPU 之间传递
    \item 利用 FlashAttention 的在线分块计算能力，边传递边计算
\end{itemize}

这与 FlashAttention 的 Tiling 思想一脉相承——你已经在作业中实现过注意力的分块计算，Ring Attention 只是将"块"分布到不同 GPU 上。

\subsection{专家并行（Expert Parallelism）}

\begin{figure}[H]
\centering
\includegraphics[width=0.95\textwidth]{slides-images/slide-051.jpg}
\caption{专家并行：MoE 模型中将不同专家分配到不同 GPU\protect\footnotemark}
\end{figure}
\footnotetext{Slides 第52页。}

专家并行用于混合专家（Mixture-of-Experts, MoE）模型。概念上类似张量并行——将一个大 MLP 拆分为多个小"专家"MLP，分布到不同 GPU。关键区别在于专家是\textbf{稀疏激活}的（只有部分专家处理每个 token），因此路由（routing）的负载均衡成为额外挑战。

\subsection{本章小结}

上下文并行解决长序列注意力的内存和计算问题，专家并行用于 MoE 架构。这些策略与前述的数据/张量/流水线并行组合使用，形成完整的 4D/5D 并行方案。

%% ========== Part 8: 组合并行策略 ==========

